%
% Template Laporan Skripsi/Thesis 
%
% @author  Andreas Febrian, Lia Sadita 
% @version 1.03
%
% Dokumen ini dibuat berdasarkan standar IEEE dalam membuat class untuk 
% LaTeX dan konfigurasi LaTeX yang digunakan Fahrurrozi Rahman ketika 
% membuat laporan skripsi. Konfigurasi yang lama telah disesuaikan dengan 
% aturan penulisan thesis yang dikeluarkan UI pada tahun 2008.
%

%
% Tipe dokumen adalah report dengan satu kolom. 
%
\documentclass[12pt, a4paper, onecolumn, oneside, final]{report}
\setlength{\parskip}{0.5em}  % Spacing after paragraphs
% Set global paragraph indentation
\setlength{\parindent}{1cm}  % Indent every paragraph by 1cm

% Load konfigurasi LaTeX untuk tipe laporan thesis
\usepackage{_internals/uithesis}

% Daftar pemenggalan suku kata dan istilah dalam LaTeX
%
% Hyphenation untuk Indonesia 
%
% @author  Andreas Febrian
% @version 1.00
% 
% Tambahkan cara pemenggalan kata-kata yang salah dipenggal secara otomatis 
% oleh LaTeX. Jika kata tersebut dapat dipenggal dengan benar, maka tidak 
% perlu ditambahkan dalam berkas ini. Tanda pemenggalan kata menggunakan 
% tanda '-'; contoh:
% menarik
%   --> pemenggalan: me-na-rik
%

\hyphenation{
    % alphabhet A
    a-na-li-sa a-tur 
    a-pli-ka-si 
    % alphabhet B
    ba-ngun-an 
    be-be-ra-pa 
    ber-ge-rak
    ber-ke-lan-jut-an 
    ber-pe-nga-ruh 
    % alphabhet C
    ca-ri
    % alphabhet D
    di-sim-pan di-pim-pin de-ngan da-e-rah di-ba-ngun da-pat di-nya-ta-kan 
    di-sim-bol-kan di-pi-lih di-li-hat de-fi-ni-si
    % alphabhet E
    e-ner-gi eks-klu-sif
    % alphabhet F
    fa-si-li-tas
    % alphabhet G
    ga-bung-an ge-rak
    % alphabhet H
    ha-lang-an
    % alphabhet I
    % alphabhet J
    % alphabhet K
    ke-hi-lang-an
    ku-ning 
    kua-li-tas ka-me-ra ke-mung-kin-an ke-se-pa-ham-an
    % alphabhet L
    ling-kung-an
    % alphabhet M
    me-neng-ah
    meng-a-tas-i me-mung-kin-kan me-nge-na-i me-ngi-rim-kan 
    meng-u-bah meng-a-dap-ta-si me-nya-ta-kan mo-di-fi-ka-si
    meng-a-tur
    % alphabhet N
    nya-ta non-eks-klu-sif
    % alphabhet O
    % alphabhet P
	pe-nye-rap-an 
	pe-ngon-trol
    pe-mo-del-an
    pe-ran  pe-ran-an-nya
    pem-ba-ngun-an pre-si-den pe-me-rin-tah prio-ri-tas peng-am-bil-an 
    peng-ga-bung-an pe-nga-was-an pe-ngem-bang-an 
    pe-nga-ruh pa-ra-lel-is-me per-hi-tung-an per-ma-sa-lah-an 
    pen-ca-ri-an peng-struk-tur-an
    % alphabhet Q
    % alphabhet R
    ran-cang-an
    % alphabhet S
    si-mu-la-si sa-ngat
    % alphabhet T
    te-ngah
    ter-da-pat
    % alphabhet U
    % alphabhet V
    % alphabhet W
    % alphabhet X
    % alphabhet Y
    % alphabhet Z
    % special
}

% Load konfigurasi khusus untuk laporan yang sedang dibuat
%-----------------------------------------------------------------------------%
% Informasi Mengenai Dokumen
%-----------------------------------------------------------------------------%
% 

% User Defined
\var{\startMagang}{25 Juni 2025}
\var{\akhirMagang}{24 Agustus 2025}
\var{\bulanSelesaiMagang}{Agustus 2025}
\var{\pembimbingHarian}{Joshevan}
\var{\pembimbingUtama}{Prof. Ray-Guang Cheng}
\var{\pembimbingUtamaTanpaProf}{Ray-Guang Cheng}
\var{\namaLab}{\textit{Broadband Multimedia Wireless Laboratory} (BMW)}
\var{\namaUniv}{\textit{National Taiwan University of Science and Technology} (NTUST)}

% Judul laporan. 
\var{\judul}{Implementasi Manajemen dan Orkestrasi Wifi pada Jaringan 5G}
% 
% Tulis kembali judul laporan, kali ini akan diubah menjadi huruf kapital
\Var{\Judul}{Implementasi Manajemen dan Orkestrasi Wifi pada Jaringan 5G}
% 
% Tulis kembali judul laporan namun dengan bahasa Ingris
\var{\judulInggris}{Implementasi of Management dan Orchestration for Wifi in 5G Network}

% 
% Tipe laporan, dapat berisi Skripsi, Tugas Akhir, Thesis, atau Disertasi
\var{\type}{Laporan Kerja Praktek}
% 
% Tulis kembali tipe laporan, kali ini akan diubah menjadi huruf kapital
\Var{\Type}{Laporan Kerja Praktek}
% 
% Tulis nama penulis 
\var{\penulis}{Edgrant Henderson Suryajaya}
% 
% Tulis kembali nama penulis, kali ini akan diubah menjadi huruf kapital
\Var{\Penulis}{Edgrant Henderson Suryajaya}
% 
% Tulis NPM penulis
\var{\npm}{2206025016}
% 
% Tuliskan Fakultas dimana penulis berada
\Var{\Fakultas}{Teknik}
\var{\fakultas}{Teknik}
% 
% Tuliskan Program Studi yang diambil penulis
\Var{\Program}{Teknik Komputer}
\Var{\program}{Teknik Komputer}
% 
% Tuliskan tahun publikasi laporan
\Var{\bulan}{Agustus}
\Var{\tahun}{2025}
% 
% Tuliskan gelar yang akan diperoleh dengan menyerahkan laporan ini
\var{\gelar}{Sarjana Teknik}
% 
% Tuliskan tanggal pengesahan laporan, waktu dimana laporan diserahkan ke 
% penguji/sekretariat
\var{\tanggalPengesahan}{23 Agustus 2025} 
% 
% Tuliskan tanggal keputusan sidang dikeluarkan dan penulis dinyatakan 
% lulus/tidak lulus
% \var{\tanggalLulus}{XX Juli 2019}
% 
% Tuliskan pembimbing 
\var{\pembimbing}{Prof. Ray}
% 
% Alias untuk memudahkan alur penulisan paa saat menulis laporan
\var{\saya}{Penulis}



%-----------------------------------------------------------------------------%
% Judul Setiap Bab
%-----------------------------------------------------------------------------%
% 
% Berikut ada judul-judul setiap bab. 
% Silahkan diubah sesuai dengan kebutuhan. 
% 
\Var{\kataPengantar}{Kata Pengantar}
\Var{\babSatu}{Pendahuluan}
\Var{\babDua}{Profil Perusahaan}
\Var{\babTiga}{Dasar Teori}
\Var{\babEmpat}{Kegiatan dan Pembahasan}
\Var{\kesimpulan}{Kesimpulan}

% Daftar istilah yang mungkin perlu ditandai 
%
% @author  Andreas Febrian
% @version 1.00
% 
% Mendaftar seluruh istilah yang mungkin akan perlu dijadikan 
% italic atau bold pada setiap kemunculannya dalam dokumen. 
% 

\var{\license}{\f{Creative Common License 1.0 Generic}}
\var{\bslash}{$\setminus$}

% Ensure the first paragraph after a section or chapter heading is indented
\usepackage{indentfirst}

% Awal bagian penulisan laporan
\begin{document}
%
% Sampul Laporan
%
% Sampul Laporan

%
% @author  unknown
% @version 1.01
% @edit by Andreas Febrian
%

\begin{titlepage}
    \begin{center}    
        \begin{figure}
            \begin{center}
                \includegraphics[width=2.5cm]{_internals/makara-warna.png}
            \end{center}
        \end{figure}    
        \vspace*{0cm}
        \bo{
        	UNIVERSITAS INDONESIA\\
        }
        
        \vspace*{1.0cm}
        % judul thesis harus dalam 14pt Times New Roman
        \large
        \bo{\Judul} \\[1.0cm]

        \vspace*{0.5 cm}    
        % harus dalam 14pt Times New Roman
        \bo{\Type}
        \normalsize

        \vspace*{6 cm}       

        % penulis dan npm
        \bo{
            \Penulis \\
            \npm \\
        }

        \vspace*{1 cm}

        % informasi mengenai fakultas dan program studi
        \setstretch{2.5} % Set line spacing
        \bo{
            DEPARTEMEN TEKNIK ELEKTRO\\
        	FAKULTAS \Fakultas \ UNIVERSITAS INDONESIA\\
        	PROGRAM STUDI \Program \\
            \vspace{2em}
        }
        \onehalfspacing
        \setstretch{1.5} % Reset line spacing

        \bo{
        	DEPOK \\
        	AGUSTUS \tahun
        }
    \end{center}
\end{titlepage}

% \var{\startMagang}{25 Juni 2025}
% \var{\akhirMagang}{24 Agustus 2025}
% \var{\bulanSelesaiMagang}{Agustus 2025}
% \var{\pembimbingHarian}{Joshevan}
% \var{\pembimbingUtama}{Prof. Ray-Guang Cheng}
% \var{\pembimbingUtamaTanpaProf}{Ray-Guang Cheng}
% \var{\namaLab}{\textit{Broadband Multimedia Wireless Laboratory} (BMW)}
% \var{\namaUniv}{\textit{National Taiwan University of Science and Technology} (NTUST)}


%
% Gunakan penomeran romawi
\pagenumbering{roman}

%
% load halaman judul dalam
\addChapter{HALAMAN JUDUL}
%
% Halaman Judul Laporan 
%
% @author  unknown
% @version 1.01
% @edit by Andreas Febrian
%

\begin{titlepage}
    \begin{center}\begin{figure}
            \begin{center}
                \includegraphics[width=2.5cm]{_internals/makara.eps}
            \end{center}
        \end{figure}    
        \vspace*{0cm}
        \bo{
        	UNIVERSITAS INDONESIA\\
        }
        
        \vspace*{1.0cm}
        % judul thesis harus dalam 14pt Times New Roman
        \bo{\Judul} \\[1.0cm]

        \vspace*{2.5 cm}    
        % harus dalam 14pt Times New Roman
        \bo{\Type} \\
        % keterangan prasyarat
        \bo{Diajukan sebagai salah satu syarat untuk memperoleh gelar \gelar}\\

        \vspace*{3 cm}       
        % penulis dan npm
        \bo{\Penulis} \\
        \bo{\npm} \\

        \vspace*{5.0cm}

        % informasi mengenai fakultas dan program studi
        \bo{
        	FAKULTAS \Fakultas\\
        	PROGRAM STUDI \Program \\
        	DEPOK \\
        	\bulan\ \tahun
        }
    \end{center}
\end{titlepage}

%
% setelah bagian ini, halaman dihitung sebagai halaman ke 2
\setcounter{page}{2}

% load halaman orisinalitas 
\addChapter{LEMBAR PERNYATAAN ORISINALITAS}
%
% Halaman Orisinalitas
%
% @author  Andreas Febrian
% @version 1.01
%

\chapter*{Halaman Pernyataan Orisinalitas}
\vspace*{2cm}

\begin{center}
	\bo{\type~ini adalah hasil karya saya sendiri, \\ 
	dan semua sumber baik yang dikutip maupun dirujuk \\
	telah saya nyatakan dengan benar.} \\
	\vspace*{2.6cm}
	
	\begin{tabular}{l c l}
	\bo{Nama} & : & \bo{\penulis} \\
	\bo{NPM} & : & \bo{\npm} \\ 
	\bo{Tanda Tangan} & : & \\
	    \includegraphics[width=0.25\linewidth]{assets/TandaTangan.png}\\

	\bo{Tanggal} & : & \bo{\tanggalPengesahan} \\	
	\end{tabular}
\end{center}

\newpage

\addChapter{LEMBAR PENGESAHAN}
%%%%%%%%%%%%%%%%%%%%%%%%%%%%%%%%%%%%%%%%%%%%%%%%%%%%%%%%
% Halaman Pengesahan Laporan Kerja Praktik
% Disesuaikan berdasarkan gambar yang diberikan.
%
% Template ini mereplikasi tata letak dokumen:
% - Judul dan informasi di tengah atas.
% - Dua kolom untuk pembimbing.
% - Satu blok di tengah bawah untuk koordinator.
% - Footer dengan nomor halaman dan nama universitas.
%%%%%%%%%%%%%%%%%%%%%%%%%%%%%%%%%%%%%%%%%%%%%%%%%%%%%%%%

% Set font size to 11pt for this page only
\fontsize{11}{16}\selectfont

% ----- Bagian Header dan Judul -----
\begin{center}
    \textbf{LEMBAR PENGESAHAN}
    
    \vspace{0.5cm}
    
    \textbf{LAPORAN KERJA PRAKTIK}
    
    \vspace{0.5cm}
    
    \textbf{\MakeUppercase{\namaUniv}} \\
    \textbf{\MakeUppercase{\namaLab}}
    
    \vspace{0.5cm}
    
    \textbf{\judul}
    
    \vspace{0.5cm}
    
    \startMagang -- \akhirMagang
\end{center}

\vspace{2cm} % Jarak vertikal sebelum blok tanda tangan

% ----- Blok Tanda Tangan (2 Kolom) -----
% Menggunakan minipage untuk membuat tata letak kolom
\noindent % Mencegah indentasi paragraf
\begin{minipage}[t]{0.5\textwidth}
    \centering
    Mengetahui, \\
    Pembimbing Utama
    
    % Sediakan ruang untuk tanda tangan. Anda bisa menyisipkan gambar di sini.
    \vspace{0.25cm}\includegraphics[width=0.5\linewidth]{assets/TandaTangan.png}\\[0.25cm]

    
    \textbf{\pembimbingUtamaTanpaProf} \\
    \textit{Head of BMW Lab NTUST}
\end{minipage}%
\begin{minipage}[t]{0.5\textwidth}
    \centering
    Menyetujui, \\
    Pembimbing Harian
    
    % Sediakan ruang untuk tanda tangan.
    \vspace{0.25cm}\includegraphics[width=0.5\linewidth]{assets/TandaTangan.png}\\[0.25cm]
    
    \textbf{\pembimbingHarian} \\
    \textit{Master Student / Mentor} \\
    \textit{BMW Lab NTUST}
\end{minipage}

\vspace{1cm} % Jarak vertikal sebelum blok tanda tangan terakhir

% ----- Blok Tanda Tangan Koordinator -----
\begin{center}
    Mengetahui, \\
    Koordinator Kerja Praktik \\
    Departemen Teknik Elektro \\
    Fakultas Teknik Universitas Indonesia
    
    % Sediakan ruang untuk tanda tangan.
    \vspace{2cm} 
    
    \textbf{Dr. Ruki Harwahyu, ST, MT, MSc.} \\
    NUP : 100120910231409891
\end{center}

% ----- Footer -----
% \vfill mendorong konten ini ke bagian bawah halaman
\vfill 


\addChapter{\kataPengantar}
%-----------------------------------------------------------------------------%
\chapter*{Kata Pengantar}
%-----------------------------------------------------------------------------%

Puji syukur saya panjatkan kepada Tuhan YME atas \ldots menyelesaikan Laporan Akhir yang berjudul “\judul”. Kegiatan ini dilakukan selama
kurang lebih 2 bulan dari \startMagang \ hingga \akhirMagang. Laporan ini ditulis
untuk memenuhi salah satu syarat untuk menyelesaikan Kerja Praktik dan juga
syarat kelulusan Fakultas Teknik Universitas Indonesia.

Selama menjalani proses \ldots dukungan dari
pihak lain \ldots terima kasih \ldots yaitu sebagai berikut:

\begin{enumerate}
	\item Tuhan YME yang telah \ldots	
        \item Bapak Dr. Ruki Harwahyu S.T., M.T., M.Sc., selaku koordinator
        program ini dari Universitas Indonesia.
	\item Ian Joseph Chandra, selaku koordinator program ini dari \namaUniv
	\item \pembimbingUtama, selaku 
        \textit{Head of }\namaLab \ sekaligus sebagai pembimbing utama dari \namaUniv
	\item \pembimbingHarian, selaku mahasiswa program master di \namaLab \ sekaligus sebagai pembimbing harian dari \namaUniv
\end{enumerate}

Adapun segala kekurangan \ldots Harapan penulis kedepannya adalah \ldots

\vspace*{1cm}
\begin{flushright}
    Taipei, 24 Agustus 2025

    \vspace{0.5cm}
    
    % --- Signature Image ---
    % Place the image directly without the 'figure' environment
    % The line break \\ ensures the name appears below it.
    \includegraphics[width=0.25\linewidth]{assets/TandaTangan.png}\\[0.5cm]

    % --- Author's Name ---
    % The author's name typically goes below the signature
    \penulis

\end{flushright}


% % load halaman pengesahan
% \addChapter{LEMBAR PERSETUJUAN}
% %
% Halaman Pengesahan
%
% @author  Andreas Febrian
% @version 1.01
%

\chapter*{HALAMAN PERSETUJUAN}

\vspace*{0.2cm}
\noindent 

\noindent
\begin{tabular}{l l p{11cm}}
	\bo{Judul}&: & \judul \\ 
	\bo{Nama}&: & \penulis \\
	\bo{NPM}&: & \npm \\
\end{tabular} \\

\vspace*{1.2cm}

\noindent Laporan \type~ini telah diperiksa dan disetujui.\\[0.3cm]
\begin{center}
\tanggalPengesahan \\[2cm]


\underline{\pembimbing}\\[0.1cm]
Pembimbing \type
\end{center}

\newpage


% \addChapter{LEMBAR PERSETUJUAN PUBLIKASI ILMIAH}
% % 
% @author  Andre Tampubolon, Andreas Febrian
% @version 1.01
% 

\chapter*{Halaman Pernyataan Persetujuan Publikasi Tugas Akhir untuk Kepentingan Akademis}

\vspace*{0.2cm}
\noindent 
Sebagai sivitas akademik Universitas Indonesia, saya yang bertanda 
tangan di bawah ini:
\vspace*{0.4cm}


\begin{tabular}{p{4.2cm} l p{6cm}}
	\bo{Nama} & : & \penulis \\ 	
	\bo{NPM} & : & \npm \\
	\bo{Program Studi} & : & \program\\	
	\bo{Fakultas} & : & \fakultas\\
	\bo{Jenis Karya} & : & \type \\
\end{tabular}

\vspace*{0.6cm}
\noindent demi pengembangan ilmu pengetahuan, menyetujui untuk memberikan 
kepada Universitas Indonesia \bo{Hak Bebas Royalti Noneksklusif 
(\textit{Non-exclusive Royalty Free Right})} atas karya ilmiah saya yang berjudul:
\begin{center}
	\judul
\end{center}
beserta perangkat yang ada (jika diperlukan). Dengan Hak Bebas Royalti 
Noneksklusif ini Universitas Indonesia berhak menyimpan, 
mengalihmedia/formatkan, mengelola dalam bentuk pangkalan data 
(\textit{database}), merawat, dan memublikasikan tugas akhir saya selama 
tetap mencantumkan nama saya sebagai penulis/pencipta dan sebagai 
pemilik Hak Cipta. \\

\noindent Demikian pernyatan ini saya buat dengan sebenarnya.

\begin{center}
	\vspace*{0.8cm}
	\begin{tabular}{rl}
		Dibuat di : & Depok \\
		Pada tanggal : & \tanggalPengesahan \\
	\end{tabular}\\

	\vspace*{0.2cm}
	Yang menyatakan \\
	\vspace*{1.1cm}
	(\penulis)
\end{center}

\newpage



 
\addChapter{ABSTRAK}
%
% Halaman Abstrak
%
% @author  Andreas Febrian
% @version 1.00
%

\chapter*{Abstrak}

\vspace*{0.2cm}
{
	\setlength{\parindent}{0pt}
	
	\begin{tabular}{@{}l l p{10cm}}
		Nama&: & \penulis \\
		Program Studi&: & \program \\
		Judul&: & \judul \\
	\end{tabular}

	\bigskip
	\bigskip

    Laporan ini membahas tentang \ldots


    \vspace{1em}
    \noindent
    \textbf{Kata Kunci:} \textit{Digital Twin}, \textit{O-RAN}, \textit{O1 Interface}, \textit{Service Management and Orchestration}, \textit{VES}, \textit{WiFi RAN}
}

\newpage


% %
% Halaman Abstract
%
% @author  Andreas Febrian
% @version 1.00
%

\chapter*{Abstract}

\vspace*{0.2cm}
{
	\setlength{\parindent}{0pt}
	
	\begin{tabular}{@{}l l p{10cm}}
		Name&: & \penulis \\
		Study Program&: & \program \\
		Title&: & \judulInggris \\
		Counsellor&: & \pembimbing \\
	\end{tabular}

	\bigskip
	\bigskip

	The focus of this study is the freshman student of Faculty of Psychology at University of
	Indonesia experience of acquiring, evaluating and using information, when they enroll in
	“Program Dasar Pendidikan Tinggi (PDPT)”. The purpose of this study is to understand
	how freshman students acquire, evaluate and use information. Knowing this will allow
	library to identify changes should be made to improve user education program at
	University of Indonesia. This research is qualitative descriptive interpretive. The data
	were collected by means of deep interview. The researcher suggests that library should
	improve the user education program and provide facilities which can help students to be
	information literate.

	\bigskip

	Key words:\\
	Information literacy, information skills, information
}

\newpage

%
% Daftar isi, gambar, dan tabel
%
\singlespacing
\phantomsection
\tableofcontents
\clearpage
\phantomsection
\listoffigures
\clearpage
\phantomsection
\listoftables
\clearpage
\phantomsection
\listoflistings
\clearpage
\onehalfspacing

%
% Gunakan penomeran Arab (1, 2, 3, ...) setelah bagian ini.
%
\pagenumbering{arabic}

%
%
%
%-----------------------------------------------------------------------------%
\chapter{\babSatu}
%-----------------------------------------------------------------------------%
Lorem ipsum dolor sit amet, consectetur adipiscing elit. Duis viverra mauris augue, et porta diam hendrerit efficitur. Nullam interdum nisl nec diam malesuada semper sed id nunc. Etiam et lacus mi. Donec eu leo ornare, sollicitudin justo in, malesuada libero. Phasellus pellentesque urna in mi luctus interdum vel eget lacus. Donec pretium erat leo, at sodales tortor imperdiet in. Nunc venenatis lacus massa, in tincidunt urna interdum sit amet. Aliquam iaculis vulputate metus a dignissim. Vivamus vulputate vitae nunc a pharetra.

Donec tincidunt nisl non elit posuere pretium. Pellentesque quis sagittis velit, quis vulputate purus. Quisque maximus justo posuere tempus dictum. Vivamus malesuada, leo eu porta elementum, purus elit tempor massa, in venenatis ligula mauris eget tortor. Phasellus venenatis nisi condimentum, tincidunt sem id, lobortis sem. Quisque volutpat, neque id ultricies finibus, ligula tortor condimentum nunc, ut placerat turpis libero eu dui. Donec faucibus blandit mauris quis posuere. Praesent a justo dapibus, blandit elit ac, pulvinar nulla. 

%-----------------------------------------------------------------------------%
\section{Latar Belakang}
%-----------------------------------------------------------------------------%
Lorem ipsum dolor sit amet, consectetur adipiscing elit. Duis viverra mauris augue, et porta diam hendrerit efficitur. Nullam interdum nisl nec diam malesuada semper sed id nunc. Etiam et lacus mi. Donec eu leo ornare, sollicitudin justo in, malesuada libero. Phasellus pellentesque urna in mi luctus interdum vel eget lacus. Donec pretium erat leo, at sodales tortor imperdiet in. Nunc venenatis lacus massa, in tincidunt urna interdum sit amet. Aliquam iaculis vulputate metus a dignissim. Vivamus vulputate vitae nunc a pharetra.

Donec tincidunt nisl non elit posuere pretium. Pellentesque quis sagittis velit, quis vulputate purus. Quisque maximus justo posuere tempus dictum. Vivamus malesuada, leo eu porta elementum, purus elit tempor massa, in venenatis ligula mauris eget tortor. Phasellus venenatis nisi condimentum, tincidunt sem id, lobortis sem. Quisque volutpat, neque id ultricies finibus, ligula tortor condimentum nunc, ut placerat turpis libero eu dui. Donec faucibus blandit mauris quis posuere. Praesent a justo dapibus, blandit elit ac, pulvinar nulla. 

%-----------------------------------------------------------------------------%
\section{Tujuan}
%-----------------------------------------------------------------------------%
Lorem ipsum dolor sit amet, consectetur adipiscing elit. Duis viverra mauris augue, et porta diam hendrerit efficitur. Nullam interdum nisl nec diam malesuada semper sed id nunc. Etiam et lacus mi. Donec eu leo ornare, sollicitudin justo in, malesuada libero. Phasellus pellentesque urna in mi luctus interdum vel eget lacus. Donec pretium erat leo, at sodales tortor imperdiet in. Nunc venenatis lacus massa, in tincidunt urna interdum sit amet. Aliquam iaculis vulputate metus a dignissim. Vivamus vulputate vitae nunc a pharetra.

Donec tincidunt nisl non elit posuere pretium. Pellentesque quis sagittis velit, quis vulputate purus. Quisque maximus justo posuere tempus dictum. Vivamus malesuada, leo eu porta elementum, purus elit tempor massa, in venenatis ligula mauris eget tortor. Phasellus venenatis nisi condimentum, tincidunt sem id, lobortis sem. Quisque volutpat, neque id ultricies finibus, ligula tortor condimentum nunc, ut placerat turpis libero eu dui. Donec faucibus blandit mauris quis posuere. Praesent a justo dapibus, blandit elit ac, pulvinar nulla. 


%-----------------------------------------------------------------------------%
\subsection{Umum}
%-----------------------------------------------------------------------------%
Lorem ipsum dolor sit amet, consectetur adipiscing elit. Duis viverra mauris augue, et porta diam hendrerit efficitur. Nullam interdum nisl nec diam malesuada semper sed id nunc. Etiam et lacus mi. Donec eu leo ornare, sollicitudin justo in, malesuada libero. Phasellus pellentesque urna in mi luctus interdum vel eget lacus. Donec pretium erat leo, at sodales tortor imperdiet in. Nunc venenatis lacus massa, in tincidunt urna interdum sit amet. Aliquam iaculis vulputate metus a dignissim. Vivamus vulputate vitae nunc a pharetra.

Donec tincidunt nisl non elit posuere pretium. Pellentesque quis sagittis velit, quis vulputate purus. Quisque maximus justo posuere tempus dictum. Vivamus malesuada, leo eu porta elementum, purus elit tempor massa, in venenatis ligula mauris eget tortor. Phasellus venenatis nisi condimentum, tincidunt sem id, lobortis sem. Quisque volutpat, neque id ultricies finibus, ligula tortor condimentum nunc, ut placerat turpis libero eu dui. Donec faucibus blandit mauris quis posuere. Praesent a justo dapibus, blandit elit ac, pulvinar nulla. 



%-----------------------------------------------------------------------------%
\subsection{Khusus}
%-----------------------------------------------------------------------------%
Lorem ipsum dolor sit amet, consectetur adipiscing elit. Duis viverra mauris augue, et porta diam hendrerit efficitur. Nullam interdum nisl nec diam malesuada semper sed id nunc. Etiam et lacus mi. Donec eu leo ornare, sollicitudin justo in, malesuada libero. Phasellus pellentesque urna in mi luctus interdum vel eget lacus. Donec pretium erat leo, at sodales tortor imperdiet in. Nunc venenatis lacus massa, in tincidunt urna interdum sit amet. Aliquam iaculis vulputate metus a dignissim. Vivamus vulputate vitae nunc a pharetra.

Donec tincidunt nisl non elit posuere pretium. Pellentesque quis sagittis velit, quis vulputate purus. Quisque maximus justo posuere tempus dictum. Vivamus malesuada, leo eu porta elementum, purus elit tempor massa, in venenatis ligula mauris eget tortor. Phasellus venenatis nisi condimentum, tincidunt sem id, lobortis sem. Quisque volutpat, neque id ultricies finibus, ligula tortor condimentum nunc, ut placerat turpis libero eu dui. Donec faucibus blandit mauris quis posuere. Praesent a justo dapibus, blandit elit ac, pulvinar nulla. 


%-----------------------------------------------------------------------------%
\section{Waktu dan Pelaksanaan Kerja Praktek}
%-----------------------------------------------------------------------------%
Lorem ipsum dolor sit amet, consectetur adipiscing elit. Duis viverra mauris augue, et porta diam hendrerit efficitur. Nullam interdum nisl nec diam malesuada semper sed id nunc. Etiam et lacus mi. Donec eu leo ornare, sollicitudin justo in, malesuada libero. Phasellus pellentesque urna in mi luctus interdum vel eget lacus. Donec pretium erat leo, at sodales tortor imperdiet in. Nunc venenatis lacus massa, in tincidunt urna interdum sit amet. Aliquam iaculis vulputate metus a dignissim. Vivamus vulputate vitae nunc a pharetra.

Donec tincidunt nisl non elit posuere pretium. Pellentesque quis sagittis velit, quis vulputate purus. Quisque maximus justo posuere tempus dictum. Vivamus malesuada, leo eu porta elementum, purus elit tempor massa, in venenatis ligula mauris eget tortor. Phasellus venenatis nisi condimentum, tincidunt sem id, lobortis sem. Quisque volutpat, neque id ultricies finibus, ligula tortor condimentum nunc, ut placerat turpis libero eu dui. Donec faucibus blandit mauris quis posuere. Praesent a justo dapibus, blandit elit ac, pulvinar nulla. 


%-----------------------------------------------------------------------------%
\section{Batasan Masalah}
%-----------------------------------------------------------------------------%
Lorem ipsum dolor sit amet, consectetur adipiscing elit. Duis viverra mauris augue, et porta diam hendrerit efficitur. Nullam interdum nisl nec diam malesuada semper sed id nunc. Etiam et lacus mi. Donec eu leo ornare, sollicitudin justo in, malesuada libero. Phasellus pellentesque urna in mi luctus interdum vel eget lacus. Donec pretium erat leo, at sodales tortor imperdiet in. Nunc venenatis lacus massa, in tincidunt urna interdum sit amet. Aliquam iaculis vulputate metus a dignissim. Vivamus vulputate vitae nunc a pharetra.

Donec tincidunt nisl non elit posuere pretium. Pellentesque quis sagittis velit, quis vulputate purus. Quisque maximus justo posuere tempus dictum. Vivamus malesuada, leo eu porta elementum, purus elit tempor massa, in venenatis ligula mauris eget tortor. Phasellus venenatis nisi condimentum, tincidunt sem id, lobortis sem. Quisque volutpat, neque id ultricies finibus, ligula tortor condimentum nunc, ut placerat turpis libero eu dui. Donec faucibus blandit mauris quis posuere. Praesent a justo dapibus, blandit elit ac, pulvinar nulla. 



%-----------------------------------------------------------------------------%
\section{Sistematika Penulisan}
%-----------------------------------------------------------------------------%
Lorem ipsum dolor sit amet, consectetur adipiscing elit. Duis viverra mauris augue, et porta diam hendrerit efficitur. Nullam interdum nisl nec diam malesuada semper sed id nunc. Etiam et lacus mi. Donec eu leo ornare, sollicitudin justo in, malesuada libero. Phasellus pellentesque urna in mi luctus interdum vel eget lacus. Donec pretium erat leo, at sodales tortor imperdiet in. Nunc venenatis lacus massa, in tincidunt urna interdum sit amet. Aliquam iaculis vulputate metus a dignissim. Vivamus vulputate vitae nunc a pharetra.

Donec tincidunt nisl non elit posuere pretium. Pellentesque quis sagittis velit, quis vulputate purus. Quisque maximus justo posuere tempus dictum. Vivamus malesuada, leo eu porta elementum, purus elit tempor massa, in venenatis ligula mauris eget tortor. Phasellus venenatis nisi condimentum, tincidunt sem id, lobortis sem. Quisque volutpat, neque id ultricies finibus, ligula tortor condimentum nunc, ut placerat turpis libero eu dui. Donec faucibus blandit mauris quis posuere. Praesent a justo dapibus, blandit elit ac, pulvinar nulla. 

%-----------------------------------------------------------------------------%

Sistematika penulisan laporan adalah sebagai berikut:
\begin{itemize}
	\item Bab 1 \babSatu \\
	\item Bab 2 \babDua \\
	\item Bab 3 \babTiga \\
	\item Bab 4 \babEmpat \\
	\item Bab 5 \kesimpulan \\
\end{itemize}

\todo{Tambahkan penjelasan singkat mengenai isi masing-masing bab.}


%-----------------------------------------------------------------------------%
\chapter{\babDua}
%-----------------------------------------------------------------------------%
\todo{tambahkan kata-kata pengantar bab 2 disini}

%-----------------------------------------------------------------------------%
\section{Profil dan Sejarah Perusahaan}
%-----------------------------------------------------------------------------%
Lorem ipsum dolor sit amet, consectetur adipiscing elit. Duis viverra mauris augue, et porta diam hendrerit efficitur. Nullam interdum nisl nec diam malesuada semper sed id nunc. Etiam et lacus mi. Donec eu leo ornare, sollicitudin justo in, malesuada libero. Phasellus pellentesque urna in mi luctus interdum vel eget lacus. Donec pretium erat leo, at sodales tortor imperdiet in. Nunc venenatis lacus massa, in tincidunt urna interdum sit amet. Aliquam iaculis vulputate metus a dignissim. Vivamus vulputate vitae nunc a pharetra.

Donec tincidunt nisl non elit posuere pretium. Pellentesque quis sagittis velit, quis vulputate purus. Quisque maximus justo posuere tempus dictum. Vivamus malesuada, leo eu porta elementum, purus elit tempor massa, in venenatis ligula mauris eget tortor. Phasellus venenatis nisi condimentum, tincidunt sem id, lobortis sem. Quisque volutpat, neque id ultricies finibus, ligula tortor condimentum nunc, ut placerat turpis libero eu dui. Donec faucibus blandit mauris quis posuere. Praesent a justo dapibus, blandit elit ac, pulvinar nulla. 

%-----------------------------------------------------------------------------%
\section{Kegiatan/Bidang Usaha Perusahaan}
%-----------------------------------------------------------------------------%
Lorem ipsum dolor sit amet, consectetur adipiscing elit. Duis viverra mauris augue, et porta diam hendrerit efficitur. Nullam interdum nisl nec diam malesuada semper sed id nunc. Etiam et lacus mi. Donec eu leo ornare, sollicitudin justo in, malesuada libero. Phasellus pellentesque urna in mi luctus interdum vel eget lacus. Donec pretium erat leo, at sodales tortor imperdiet in. Nunc venenatis lacus massa, in tincidunt urna interdum sit amet. Aliquam iaculis vulputate metus a dignissim. Vivamus vulputate vitae nunc a pharetra.

Donec tincidunt nisl non elit posuere pretium. Pellentesque quis sagittis velit, quis vulputate purus. Quisque maximus justo posuere tempus dictum. Vivamus malesuada, leo eu porta elementum, purus elit tempor massa, in venenatis ligula mauris eget tortor. Phasellus venenatis nisi condimentum, tincidunt sem id, lobortis sem. Quisque volutpat, neque id ultricies finibus, ligula tortor condimentum nunc, ut placerat turpis libero eu dui. Donec faucibus blandit mauris quis posuere. Praesent a justo dapibus, blandit elit ac, pulvinar nulla. 

%-----------------------------------------------------------------------------%
\section{Lokasi Perusahaan}
%-----------------------------------------------------------------------------%
Lorem ipsum dolor sit amet, consectetur adipiscing elit. Duis viverra mauris augue, et porta diam hendrerit efficitur. Nullam interdum nisl nec diam malesuada semper sed id nunc. Etiam et lacus mi. Donec eu leo ornare, sollicitudin justo in, malesuada libero. Phasellus pellentesque urna in mi luctus interdum vel eget lacus. Donec pretium erat leo, at sodales tortor imperdiet in. Nunc venenatis lacus massa, in tincidunt urna interdum sit amet. Aliquam iaculis vulputate metus a dignissim. Vivamus vulputate vitae nunc a pharetra.

Donec tincidunt nisl non elit posuere pretium. Pellentesque quis sagittis velit, quis vulputate purus. Quisque maximus justo posuere tempus dictum. Vivamus malesuada, leo eu porta elementum, purus elit tempor massa, in venenatis ligula mauris eget tortor. Phasellus venenatis nisi condimentum, tincidunt sem id, lobortis sem. Quisque volutpat, neque id ultricies finibus, ligula tortor condimentum nunc, ut placerat turpis libero eu dui. Donec faucibus blandit mauris quis posuere. Praesent a justo dapibus, blandit elit ac, pulvinar nulla. 

%-----------------------------------------------------------------------------%
\section{Struktur Organisasi Perusahaan}
%-----------------------------------------------------------------------------%
Lorem ipsum dolor sit amet, consectetur adipiscing elit. Duis viverra mauris augue, et porta diam hendrerit efficitur. Nullam interdum nisl nec diam malesuada semper sed id nunc. Etiam et lacus mi. Donec eu leo ornare, sollicitudin justo in, malesuada libero. Phasellus pellentesque urna in mi luctus interdum vel eget lacus. Donec pretium erat leo, at sodales tortor imperdiet in. Nunc venenatis lacus massa, in tincidunt urna interdum sit amet. Aliquam iaculis vulputate metus a dignissim. Vivamus vulputate vitae nunc a pharetra.

Donec tincidunt nisl non elit posuere pretium. Pellentesque quis sagittis velit, quis vulputate purus. Quisque maximus justo posuere tempus dictum. Vivamus malesuada, leo eu porta elementum, purus elit tempor massa, in venenatis ligula mauris eget tortor. Phasellus venenatis nisi condimentum, tincidunt sem id, lobortis sem. Quisque volutpat, neque id ultricies finibus, ligula tortor condimentum nunc, ut placerat turpis libero eu dui. Donec faucibus blandit mauris quis posuere. Praesent a justo dapibus, blandit elit ac, pulvinar nulla. 

%-----------------------------------------------------------------------------%
\section{Profil Divisi Real Time Laboratory}
%-----------------------------------------------------------------------------%
Lorem ipsum dolor sit amet, consectetur adipiscing elit. Duis viverra mauris augue, et porta diam hendrerit efficitur. Nullam interdum nisl nec diam malesuada semper sed id nunc. Etiam et lacus mi. Donec eu leo ornare, sollicitudin justo in, malesuada libero. Phasellus pellentesque urna in mi luctus interdum vel eget lacus. Donec pretium erat leo, at sodales tortor imperdiet in. Nunc venenatis lacus massa, in tincidunt urna interdum sit amet. Aliquam iaculis vulputate metus a dignissim. Vivamus vulputate vitae nunc a pharetra.

Donec tincidunt nisl non elit posuere pretium. Pellentesque quis sagittis velit, quis vulputate purus. Quisque maximus justo posuere tempus dictum. Vivamus malesuada, leo eu porta elementum, purus elit tempor massa, in venenatis ligula mauris eget tortor. Phasellus venenatis nisi condimentum, tincidunt sem id, lobortis sem. Quisque volutpat, neque id ultricies finibus, ligula tortor condimentum nunc, ut placerat turpis libero eu dui. Donec faucibus blandit mauris quis posuere. Praesent a justo dapibus, blandit elit ac, pulvinar nulla. 

%-----------------------------------------------------------------------------%
% \section{\latex~Secara Singkat}
% %-----------------------------------------------------------------------------%
% Definisi dari LaTeX \citep{lankton2008introduction} adalah: \\ 
% \begin{tabular}{| p{13cm} |}
% 	\hline 
% 	\\
% 	LaTeX is a family of programs designed to produce publication-quality 
% 	typeset documents. It is particularly strong when working with 
% 	mathematical symbols. \\	
% 	The history of LaTeX begins with a program called TEX. In 1978, a 
% 	computer scientist by the name of Donald Knuth grew frustrated with the 
% 	mistakes that his publishers made in typesetting his work. He decided 
% 	to create a typesetting program that everyone could easily use to 
% 	typeset documents, particularly those that include formulae, and made 
% 	it freely available. The result is TEX. \\	
% 	Knuth's product is an immensely powerful program, but one that does 
% 	focus very much on small details. A mathematician and computer 
% 	scientist by the name of Leslie Lamport wrote a variant of TEX called 
% 	LaTeX that focuses on document structure rather than such details. \\
% 	\\
% 	\hline
% \end{tabular}

% \vspace*{0.8cm}

% Contoh sitasi lainnya menggunakan \verb|\citep| adalah saat kita mau mensitasi pekerjaan tentang \textit{machine learning} \citep{chin2000learning} dan \textit{dynamic programming} \citep{barto1995learning}. \\

% Dokumen \latex~sangat mudah, seperti halnya membuat dokumen teks biasa. Ada 
% beberapa perintah yang diawali dengan tanda '\bslash'. 
% Seperti perintah \bslash\bslash~yang digunakan untuk memberi baris baru. 
% Perintah tersebut juga sama dengan perintah \bslash newline. 
% Pada bagian ini akan sedikit dijelaskan cara manipulasi teks dan 
% perintah-perintah \latex~yang mungkin akan sering digunakan. 
% Jika ingin belajar hal-hal dasar mengenai \latex, silahkan kunjungi: 

% \begin{itemize}
% 	\item \url{http://frodo.elon.edu/tutorial/tutorial/}, atau
% 	\item \url{http://www.maths.tcd.ie/~dwilkins/LaTeXPrimer/}
% \end{itemize}


% %-----------------------------------------------------------------------------%
% \section{\latex~Kompiler dan IDE}
% %-----------------------------------------------------------------------------%
% Agar dapat menggunakan \latex~(pada konteks hanya sebagai pengguna), Anda 
% tidak perlu banyak tahu mengenai hal-hal didalamnya. 
% Seperti halnya pembuatan dokumen secara visual (contohnya Open Office (OO) 
% Writer), Anda dapat menggunakan \latex~dengan cara yang sama. 
% Orang-orang yang menggunakan \latex~relatif lebih teliti dan terstruktur 
% mengenai cara penulisan yang dia gunakan, \latex~memaksa Anda untuk seperti 
% itu.  

% Kembali pada bahasan utama, untuk mencoba \latex~Anda cukup mendownload 
% kompiler dan IDE. Saya menyarankan menggunakan Texlive dan Texmaker. 
% Texlive dapat didownload dari \url{http://www.tug.org/texlive/}. 
% Sedangkan Texmaker dapat didownload dari 
% \url{http://www.xm1math.net/texmaker/}. 
% Untuk pertama kali, coba buka berkas laporan-kp.tex dalam template yang Anda miliki 
% pada Texmaker. 
% Dokumen ini adalah dokumen utama. 
% Tekan F6 (PDFLaTeX) dan Texmaker akan mengkompilasi berkas tersebut menjadi 
% berkas PDF. 
% Jika tidak bisa, pastikan Anda sudah menginstall Texlive. 
% Buka berkas tersebut dengan menekan F7. 
% Hasilnya adalah sebuah dokumen yang sama seperti dokumen yang Anda baca saat 
% ini. 


% %-----------------------------------------------------------------------------%
% \section{Bold, Italic, dan Underline}
% %-----------------------------------------------------------------------------%
% Hal pertama yang mungkin ditanyakan adalah bagaimana membuat huruf tercetak 
% tebal, miring, atau memiliki garis bawah. 
% Pada Texmaker, Anda bisa melakukan hal ini seperti halnya saat mengubah dokumen 
% dengan OO Writer. 
% Namun jika tetap masih tertarik dengan cara lain, ini dia: 

% \begin{itemize}
% 	\item \bo{Bold} \\
% 		Gunakan perintah \bslash textbf$\lbrace\rbrace$ atau 
% 		\bslash bo$\lbrace\rbrace$. 
% 	\item \f{Italic} \\
% 		Gunakan perintah \bslash textit$\lbrace\rbrace$ atau 
% 		\bslash f$\lbrace\rbrace$. 
% 	\item \underline{Underline} \\
% 		Gunakan perintah \bslash underline$\lbrace\rbrace$.
% 	\item $\overline{Overline}$ \\
% 		Gunakan perintah \bslash overline. 
% 	\item $^{superscript}$ \\
% 		Gunakan perintah \bslash $\lbrace\rbrace$. 
% 	\item $_{subscript}$ \\
% 		Gunakan perintah \bslash \_$\lbrace\rbrace$. 
% \end{itemize}

% Perintah \bslash f dan \bslash bo hanya dapat digunakan jika package 
% uithesis digunakan. 


% %-----------------------------------------------------------------------------%
% \section{Memasukan Gambar}
% %-----------------------------------------------------------------------------%
% Setiap gambar dapat diberikan caption dan diberikan label. Label dapat 
% digunakan untuk menunjuk gambar tertentu. 
% Jika posisi gambar berubah, maka nomor gambar juga akan diubah secara 
% otomatis. 
% Begitu juga dengan seluruh referensi yang menunjuk pada gambar tersebut. 
% Contoh sederhana adalah \pic~\ref{fig:testGambar}. 
% Silahkan lihat code \latex~dengan nama bab2.tex untuk melihat kode lengkapnya. 
% Harap diingat bahwa caption untuk gambar selalu terletak dibawah gambar. 

% \begin{figure}
% 	\centering
% 	\includegraphics[width=0.50\textwidth]
% 		{assets/pics/creative_common.png}
% 	\caption{\license.}
% 	\label{fig:testGambar}
% \end{figure}


% %-----------------------------------------------------------------------------%
% \section{Membuat Tabel}
% %-----------------------------------------------------------------------------%
% Seperti pada gambar, tabel juga dapat diberi label dan caption. 
% Caption pada tabel terletak pada bagian atas tabel. 
% Contoh tabel sederhana dapat dilihat pada \tab~\ref{tab:tab1}.

% \begin{table}
% 	\centering
% 	\caption{Contoh Tabel}
% 	\label{tab:tab1}
% 	\begin{tabular}{| l | c r |}
% 		\hline
% 		& kol 1 & kol 2 \\ 
% 		\hline
% 		baris 1 & 1 & 2 \\
% 		baris 2 & 3 & 4 \\
% 		baris 3 & 5 & 6 \\
% 		jumlah  & 9 & 12 \\
% 		\hline
% 	\end{tabular}
% \end{table}

% Ada jenis tabel lain yang dapat dibuat dengan \latex~berikut 
% beberapa diantaranya. 
% Contoh-contoh ini bersumber dari 
% \url{http://en.wikibooks.org/wiki/LaTeX/Tables}

% \begin{table}
% 	\centering
% 	\caption{An Example of Rows Spanning Multiple Columns}
% 	\label{row.spanning}
% 	\begin{tabular}{|l|l|*{6}{c|}}
%   		\hline % create horizontal line
%   		No & Name & \multicolumn{3}{|c|}{Week 1} & \multicolumn{3}{|c|}{Week 2} \\
%   		\cline{3-8} % create line from 3rd column till 8th column
%   		& & A & B & C & A & B & C\\
%   		\hline
%   		1 & Lala & 1 & 2 & 3 & 4 & 5 & 6\\
%   		2 & Lili & 1 & 2 & 3 & 4 & 5 & 6\\
%   		3 & Lulu & 1 & 2 & 3 & 4 & 5 & 6\\
%   		\hline
% 	\end{tabular}
% \end{table}

% \begin{table}
% 	\centering
% 	\caption{An Example of Columns Spanning Multiple Rows}
% 	\label{column.spanning}
% 	\begin{tabular}{|l|c|l|}
% 		\hline
% 		Percobaan & Iterasi & Waktu \\
% 		\hline
% 		Pertama & 1 & 0.1 sec \\ \hline
% 		\multirow{2}{*}{Kedua} & 1 & 0.1 sec \\
%  		& 3 & 0.15 sec \\ 
%  		\hline
% 		\multirow{3}{*}{Ketiga} & 1 & 0.09 sec \\
%  		& 2 & 0.16 sec \\
%  		& 3 & 0.21 sec \\ 
%  		\hline
% 	\end{tabular}
% \end{table}

% \begin{table}
% 	\centering
% 	\caption{An Example of Spanning in Both Directions Simultaneously}
% 	\label{mix.spanning}
% 	\begin{tabular}{cc|c|c|c|c|}
% 		\cline{3-6}
% 		& & \multicolumn{4}{|c|}{Title} \\ \cline{3-6}
% 		& & A & B & C & D \\ \hline
% 		\multicolumn{1}{|c|}{\multirow{2}{*}{Type}} &
% 		\multicolumn{1}{|c|}{X} & 1 & 2 & 3 & 4\\ \cline{2-6}
% 		\multicolumn{1}{|c|}{}                        &
% 		\multicolumn{1}{|c|}{Y} & 0.5 & 1.0 & 1.5 & 2.0\\ \cline{1-6}
% 		\multicolumn{1}{|c|}{\multirow{2}{*}{Resource}} &
% 		\multicolumn{1}{|c|}{I} & 10 & 20 & 30 & 40\\ \cline{2-6}
% 		\multicolumn{1}{|c|}{}                        &
% 		\multicolumn{1}{|c|}{J} & 5 & 10 & 15 & 20\\ \cline{1-6}
% 	\end{tabular}
% \end{table}


%-----------------------------------------------------------------------------%
\chapter{\babTiga}
%-----------------------------------------------------------------------------%
\todo{tambahkan kata-kata pengantar bab 1 disini}


%-----------------------------------------------------------------------------%
\section{Satu Persamaan}
%-----------------------------------------------------------------------------%

\noindent \begin{align}\label{eq:garis}
	\cfrac{y - y_{1}}{y_{2} - y_{1}} = 
	\cfrac{x - x_{1}}{x_{2} - x_{1}}
\end{align}

\equ~\ref{eq:garis} diatas adalah persamaan garis. 
\equ~\ref{eq:garis} dan \ref{eq:bola} sama-sama dibuat dengan perintah \bslash
align. 
Perintah ini juga dapat digunakan untuk menulis lebih dari satu persamaan. 

\noindent \begin{align}\label{eq:bola}
	\underbrace{|\overline{ab}|}_{\text{pada bola $|\overline{ab}| = r$}} 
		= \sqrt[2]{(x_{b} - x_{a})^{2} + (y_{b} - y_{a})^{2} + 
				\vert\vert(z_{b} - z_{a})^{2}}
\end{align}

%-----------------------------------------------------------------------------%
\section{Lebih dari Satu Persamaan}
\label{sec:multiEqu}
%-----------------------------------------------------------------------------%
\noindent \begin{align}\label{eq:matriks}	
	|\overline{a} * \overline{b}| &= |\overline{a}| |\overline{b}| \sin\theta 
		\\[0.2cm]
	\overline{a} * \overline{b} &=  
		\begin{array}{| c c c |}
			\hat{i} & x_{1} & x_{2} \\
			\hat{j} & y_{1} & y_{2} \\
			\hat{k} & z_{1} & z_{2} \\
		\end{array} \nonumber \\[0.2cm]
	&= \hat{i} \,
		\begin{array}{ | c c | }
			y_{1} & y_{2} \\
			z_{1} & z_{2} \\
		\end{array} 
	   + \hat{j} \,
		\begin{array}{ | c c | }
			z_{1} & z_{2} \\
			x_{1} & x_{2} \\
		\end{array} 
	   + \hat{k} \,	
		\begin{array}{ | c c | }
			x_{1} & x_{2} \\
			y_{1} & y_{2} \\
		\end{array}
		\nonumber
\end{align}

Pada \equ~\ref{eq:matriks} dapat dilihat beberapa baris menjadi satu bagian 
dari \equ~\ref{eq:matriks}. 
Sedangkan dibawah ini dapat dilihat bahwa dengan cara yang sama, \equ~
\ref{eq:gabungan1}, \ref{eq:gabungan2}, dan \ref{eq:gabungan3} memiliki nomor 
persamaannya masing-masing. 

\noindent \begin{align}\label{eq:gabungan1}	
	\int_{a}^{b} f(x)\, dx + \int_{b}^{c} f(x) \, dx = \int_{a}^{c} f(x) \, dx
		\\\label{eq:gabungan2}
	\lim_{x \to \infty} \frac{f(x)}{g(x)} = 0 \hspace{1cm} 
		\text{jika pangkat $f(x)$ $<$ pangkat $g(x)$} \\\label{eq:gabungan3}
	a^{m^{a \, ^{n}\log b }} = b^{\frac{m}{n}}
\end{align}


%-----------------------------------------------------------------------------%
\chapter{\babEmpat}
%-----------------------------------------------------------------------------%
\todo{tambahkan kata-kata pengantar bab 1 disini}

%-----------------------------------------------------------------------------%
\section{thesis.tex}
%-----------------------------------------------------------------------------%
Berkas ini berisi seluruh berkas Latex yang dibaca, jadi bisa dikatakan sebagai 
berkas utama. Dari berkas ini kita dapat mengatur bab apa saja yang ingin 
kita tampilkan dalam dokumen.


%-----------------------------------------------------------------------------%
\section{laporan\_setting.tex}
%-----------------------------------------------------------------------------%
Berkas ini berguna untuk mempermudah pembuatan beberapa template standar. 
Anda diminta untuk menuliskan judul laporan, nama, npm, dan hal-hal lain yang 
dibutuhkan untuk pembuatan template. 


%-----------------------------------------------------------------------------%
\section{istilah.tex}
%-----------------------------------------------------------------------------%
Berkas istilah digunakan untuk mencatat istilah-istilah yang digunakan. 
Fungsinya hanya untuk memudahkan penulisan.
Pada beberapa kasus, ada kata-kata yang harus selalu muncul dengan tercetak 
miring atau tercetak tebal. 
Dengan menjadikan kata-kata tersebut sebagai sebuah perintah \latex~tentu akan 
mempercepat dan mempermudah pengerjaan laporan. 


%-----------------------------------------------------------------------------%
\section{hype.indonesia.tex}
%-----------------------------------------------------------------------------%
Berkas ini berisi cara pemenggalan beberapa kata dalam bahasa Indonesia. 
\latex~memiliki algoritma untuk memenggal kata-kata sendiri, namun untuk 
beberapa kasus algoritma ini memenggal dengan cara yang salah. 
Untuk memperbaiki pemenggalan yang salah inilah cara pemenggalan yang benar 
ditulis dalam berkas hype.indonesia.tex.


%-----------------------------------------------------------------------------%
\section{pustaka.tex}
%-----------------------------------------------------------------------------%
Berkas pustaka.tex berisi seluruh daftar referensi yang digunakan dalam 
laporan. 
Anda bisa membuat model daftar referensi lain dengan menggunakan bibtex.
Untuk mempelajari bibtex lebih lanjut, silahkan buka 
\url{http://www.bibtex.org/Format}. 
Untuk merujuk pada salah satu referensi yang ada, gunakan perintah \bslash 
cite, e.g. \bslash cite\{lankton2008introduction\} yang akan akan memunculkan 
\cite{lankton2008introduction}


%-----------------------------------------------------------------------------%
\section{bab[1 - 6].tex}
%-----------------------------------------------------------------------------%
Berkas ini berisi isi laporan yang Anda tulis. 
Setiap nama berkas e.g. bab1.tex merepresentasikan bab dimana tulisan tersebut 
akan muncul. 
Sebagai contoh, kode dimana tulisan ini dibaut berada dalam berkas dengan nama 
bab4.tex. 
Ada enam buah berkas yang telah disiapkan untuk mengakomodir enam bab dari 
laporan Anda, diluar bab kesimpulan dan saran. 
Jika Anda tidak membutuhkan sebanyak itu, silahkan hapus kode dalam berkas 
thesis.tex yang memasukan berkas \latex~yang tidak dibutuhkan;  contohnya 
perintah \bslash include\{bab6.tex\} merupakan kode untuk memasukan berkas 
bab6.tex kedalam laporan.

%-----------------------------------------------------------------------------%
\section{Penulisan \textit{code} atau \textit{pseudocode} program}
%-----------------------------------------------------------------------------%

\subsection{\textit{Inline}}

Dengan perintah \verb|\verb|: \verb|System.out.println("Hello, World");| \\
Dengan perintah \textit{custom} \verb|\code|: \code{System.out.println("Hello, World"); }
Dengan perintah \verb|\mintinline|: \mintinline{java}{System.out.println("Hello, World"); }

\subsection{\textit{Multiline}}

Dengan perintah \verb|verbatim|: 

\begin{verbatim}	
public class HelloWorld {
    public static void main(String[] args) {
        // Prints "Hello, World" to the terminal window.
        System.out.println("Hello, World");
    }
}
\end{verbatim}

Dengan perintah \verb|minted|: Kode \ref{code:hw:minted}
\begin{listing}[H]
    \begin{minted}{python}
def binary_accuracy(y_true, y_pred):
    return K.mean(K.equal(y_true, K.round(y_pred)), axis=-1)


def categorical_accuracy(y_true, y_pred):
    return K.cast(K.equal(K.argmax(y_true, axis=-1),
                          K.argmax(y_pred, axis=-1)),
                  K.floatx())


def sparse_categorical_accuracy(y_true, y_pred):
    # reshape in case it's in shape (num_samples, 1) instead of (num_samples,)
    if K.ndim(y_true) == K.ndim(y_pred):
        y_true = K.squeeze(y_true, -1)
    # convert dense predictions to labels
    y_pred_labels = K.argmax(y_pred, axis=-1)
    y_pred_labels = K.cast(y_pred_labels, K.floatx())
    return K.cast(K.equal(y_true, y_pred_labels), K.floatx())


def top_k_categorical_accuracy(y_true, y_pred, k=5):
    return K.mean(K.in_top_k(y_pred, K.argmax(y_true, axis=-1), k), axis=-1)


def sparse_top_k_categorical_accuracy(y_true, y_pred, k=5):
    # If the shape of y_true is (num_samples, 1), flatten to (num_samples,)
    return K.mean(K.in_top_k(y_pred, K.cast(K.flatten(y_true), 'int32'), k),
                  axis=-1)
    \end{minted}
    \caption{An excerpt from keras: \url{https://github.com/keras-team/keras/blob/master/keras/metrics.py}}
    \label{code:hw:minted}
\end{listing}

Konfigurasi tampilan bisa dilakukan di \verb|uithesis.sty| dengan referensi dokumentasi di \url{https://github.com/gpoore/minted/blob/master/source/minted.pdf}
%---------------------------------------------------------------
\chapter{\kesimpulan}
%---------------------------------------------------------------
\todo{Tambahkan kesimpulan dan saran terkait dengan perkerjaan 
	yang dilakukan.}


%---------------------------------------------------------------
\section{Kesimpulan}
%---------------------------------------------------------------


%---------------------------------------------------------------
\section{Saran}
%---------------------------------------------------------------


%
% Daftar Pustaka
%\input{pustaka}

% Alternatif manajemen daftar pustaka dengan \bibliography
\bibliographystyle{apalike}
\bibliography{pustaka}

%
% Lampiran 
%
% \begin{appendix}
% 	%
% @author  Andreas Febrian
% @version 1.00 
% 
% Hanya sebuah pembatas bertuliskan LAMPIRAN ditengah halaman. 
% 

\begin{titlepage}
	\centering 
	\vspace*{6cm}
	\noindent \Huge{LAMPIRAN}
	\addChapter{LAMPIRAN}
\end{titlepage}
% 	\setcounter{page}{2}\textbf{}
% 	%-----------------------------------------------------------------------------%
\addChapter{Lampiran 1}
\chapter*{Lampiran 1}
%-----------------------------------------------------------------------------%
% \end{appendix}

\end{document}